\section{Setting: typed transition data and assessment operators}\label{sec:setting}

\subsection{Typed transition data and exact-tube semantics}

A \emph{typed transition datum} consists of a state space $Z$ whose coordinates carry types (a conserved stock with a unit, a service margin, an institutional label), a finite \emph{menu} $\mathcal A$ of plans, a \emph{disturbance class} $D$, and, for each plan $a$ and disturbance $d$, a \emph{tube} $\Tube(a,d)\subseteq Z$ -- the full set of states visited during one review interval -- together with its \emph{endpoint set} $\End(a,d)\subseteq\Tube(a,d)$ and its \emph{successor set} $\Succ(a,d)$, the states reached after any end-of-interval reset. Among the coordinates are $n$ \emph{typed floors} $s=(s_1,\dots,s_n)$, each normalised so that $s_i\ge0$ is the binding constraint and each carrying the provenance of its floor (physical, contractual or normative). Typing records which quantities may be added: conservation statements are per type, and the framework does not itself authorise adding biomass, money and biodiversity into one conserved scalar; the scalarization weights studied below are precisely such a bridge, declared as an assessment convention rather than a physical identity.

Two constraint sets complete the datum: a \emph{physical safe set} $S^{\mathrm{phys}}$ and a \emph{physical destination set} $G^{\mathrm{phys}}$, which encode constraints that no assessment doctrine disputes (for instance that a reserve stock cannot be negative). The \emph{typed} sets add the floors,
\[
S=S^{\mathrm{phys}}\cap\{s_i\ge0,\ i=1,\dots,n\},\qquad G=G^{\mathrm{phys}}\cap\{s\ge0\},
\]
and the \emph{scalarized} sets at a weight $w$ add only the aggregate floor,
\[
S^w=S^{\mathrm{phys}}\cap\{w\cdot s\ge0\},\qquad G^w=G^{\mathrm{phys}}\cap\{w\cdot s\ge0\}.
\]
Throughout, $w$ ranges over the full nonnegative cone $\Wp=\Rp\setminus\{0\}$ unless a subfamily $W\subseteq\Wp$ is named. The semantics is \emph{exact-tube}: a plan is admissible only if every visited state, not merely the endpoint, lies in the relevant safe set. The disturbance quantifier is always innermost -- admissibility requires safety for every $d\in D$ -- and the separation studied here concerns a second quantifier, over weights, placed outside it.

\subsection{Assessment operators}

For a state $z$, four operators return the set of plans that an assessment doctrine accepts at $z$:
\begin{align}
\Etyp(z)&=\{a:\ \forall d\in D,\ \Tube(a,d)\subseteq S,\ \Succ(a,d)\subseteq G\}, \label{eq:Etyp}\\
E_w(z)&=\{a:\ \forall d\in D,\ \Tube(a,d)\subseteq S^w,\ \Succ(a,d)\subseteq G^w\},\quad w\in\Wp, \label{eq:Ew}\\
\Ephys(z)&=\{a:\ \forall d\in D,\ \Tube(a,d)\subseteq S^{\mathrm{phys}},\ \Succ(a,d)\subseteq G^{\mathrm{phys}}\}, \label{eq:Ephys}\\
\Eend(z)&=\{a:\ \forall d\in D,\ \End(a,d)\subseteq S^{\mathrm{phys}},\ \Succ(a,d)\subseteq G^{\mathrm{phys}}\}. \label{eq:Eend}
\end{align}
$\Etyp$ is the \emph{noncompensatory typed} operator (each floor separately binding along the path); $E_w$ the \emph{scalarized} operator at weight $w$ (the dashboard reading $w\cdot s$ must stay nonnegative along the path); $\Ephys$ the \emph{physical} operator (no floors at all); and $\Eend$ the \emph{endpoint} operator, which represents any assessment architecture that evaluates constraints only at sampled or audited endpoints. The two axes along which they differ -- what is constrained and where it is evaluated -- are set out in Table~\ref{tab:semantics}; the typed-endpoint operator $E_{\mathrm{end,typ}}$, which applies the floors at endpoints only, is the remaining cell and is used once, in Section~\ref{sec:design}.

\begin{table}[!htbp]
\centering\small
\caption{Evaluation semantics: constraint structure (rows) by evaluation locus (columns).}
\label{tab:semantics}
\begin{tabularx}{\textwidth}{@{}>{\raggedright\arraybackslash}p{0.26\textwidth}XX@{}}
\toprule
 & Evaluated on the full tube & Evaluated at endpoints only \\
\midrule
Typed floors & $\Etyp$ -- typed path safety & $E_{\mathrm{end,typ}}$ -- typed snapshots \\
Scalarized floor & $E_w$ -- scalarized path safety & (not isolated here) \\
Physical constraints only & $\Ephys$ -- physical viability & $\Eend$ -- audited physical snapshots \\
\bottomrule
\end{tabularx}
\end{table}

\subsection{Accepted-state sets and the central inclusion}

For any operator $E$ write $\mathcal V[E]=\{z:\ E(z)\ne\varnothing\}$ for its \emph{accepted-state set}, the states at which the doctrine admits at least one plan. Thus $\Vtyp=\mathcal V[\Etyp]$, $\mathcal V_w=\mathcal V[E_w]$, $\Vphys=\mathcal V[\Ephys]$, $\Vend=\mathcal V[\Eend]$, and the \emph{compensatory accepted set} is the per-weight intersection
\begin{equation}
\Vweak=\bigcap_{w\in\Wp}\mathcal V_w=\{z:\ \forall w\ \exists a_w\in E_w(z)\}.
\label{eq:Vweak}
\end{equation}
The object of the paper is the inclusion
\begin{equation}
\underbrace{\{z:\ \textstyle\bigcap_w E_w(z)\neq\varnothing\}}_{\text{common plan}}\ \subseteq\ \underbrace{\bigcap_w\{z:\ E_w(z)\neq\varnothing\}}_{\text{per weight}\ =\ \Vweak},
\label{eq:central}
\end{equation}
whose left-hand side will turn out to be exactly $\Vtyp$ (Proposition~\ref{prop:hierarchy}(ii)), and the set $\FP=\Vweak\setminus\Vtyp$ on which it is strict -- the \emph{acceptance gap}. Table~\ref{tab:notation} collects the symbols.

\begin{table}[!htbp]
\centering\small
\caption{Notation.}
\label{tab:notation}
\begin{tabularx}{\textwidth}{@{}>{\raggedright\arraybackslash}p{0.34\textwidth}X@{}}
\toprule
Symbol & Meaning \\
\midrule
$z$; $a\in\mathcal A$; $d\in D$ & state; plan (action) from the finite menu; disturbance from the declared class \\
$s=(s_1,\dots,s_n)$; $x$ & typed floors, normalised so that $s_i\ge0$ binds; reserve (bridge) stock on the witness \\
$\Wp=\Rp\setminus\{0\}$; $W\subseteq\Wp$; $r=w_2/w_1$ & full weight cone; a weight subfamily; weight ratio on the two-floor witness \\
$S,G$; $S^w,G^w$; $S^{\mathrm{phys}},G^{\mathrm{phys}}$ & typed, scalarized and physical safe/destination sets \\
$\Tube,\End,\Succ$ & visited set, endpoint set and successor set of a plan under a disturbance \\
$\Etyp$, $E_w$, $\Ephys$, $\Eend$ & assessment operators \eqref{eq:Etyp}--\eqref{eq:Eend}; $\mathcal V[E]=\{z:E(z)\ne\varnothing\}$ \\
$\Vtyp$, $\mathcal V_w$, $\Vweak$, $\Vphys$, $\Vend$ & accepted-state sets; $\Vweak=\bigcap_{w\in\Wp}\mathcal V_w$ \\
$\FP=\Vweak\setminus\Vtyp$ & the acceptance gap \\
$d_a$; $\mathcal D=\{d_a\}$ & required-margin vector of plan $a$ (Theorem~\ref{thm:menu}) \\
$\Q,\ R,\ I$; $\rho_1,\rho_2$ & discrepancy region, rescue set, impossibility region; licensing thresholds (Theorem~\ref{thm:witness}) \\
$e=(\tfrac14,\tfrac14)$; $c=1$; $\kappa,\kappa^*$ & destination gain and rescue cost of the witness; reserve increment and rescue threshold \\
$\theta$; $\sigma=1/(1-\theta)$; $\lambda_i=1+s_i$ & CES exponent, elasticity of substitution, floor-referenced index (Theorem~\ref{thm:ladder}) \\
\bottomrule
\end{tabularx}
\end{table}

